\begingroup
\makeatletter\tagsleft@true\let\veqno\@@leqno\makeatother
\ArithmeticHead{351}{R}{{\scriptsize ZUR THEORIE DER REIHE $F(\alpha,\beta,\gamma,x)$.\quad III.}}
\NumberedItem{8}
\noindent{\small[S. 67]}\par
\indent Zu den oben gelieferten Untersuchungen können wir noch folgende wichtige Nachträge liefern.
\begin{equation}
(1+x)^\alpha F(\alpha,\alpha+1-\gamma,\gamma,x)=F\left(\tfrac12\alpha,\tfrac12\alpha+\tfrac12,\gamma,\frac{4x}{(1+x)^2}\right),\tag*{I.}
\end{equation}
\begin{equation}
(1+x)^{2\alpha}F(\alpha,\alpha+1-\gamma,\gamma,xx)=F\left(\alpha,\gamma-\tfrac12,2\gamma-1,\frac{4x}{(1+x)^2}\right),\tag*{II.}
\end{equation}
\begin{equation}
\begin{aligned}
&F(\alpha,\beta,\gamma,1)=\frac{(\gamma-\alpha)(\gamma-\beta)}{\gamma(\gamma-\alpha-\beta)}\cdot\frac{(\gamma+1-\alpha)(\gamma+1-\beta)}{(\gamma+1)(\gamma+1-\alpha-\beta)}\\[.7em]
&\qquad\cdot\frac{(\gamma+2-\alpha)(\gamma+2-\beta)}{(\gamma+2)(\gamma+2-\alpha-\beta)}\ \mathrm{etc.}\\[.7em]
&\qquad=\frac{\Pi(\gamma-1)\Pi(\gamma-1-\alpha-\beta)}{\Pi(\gamma-1-\alpha)\Pi(\gamma-1-\beta)},
\end{aligned}\tag*{III.}
\end{equation}
\noindent wenn man durch $\Pi z$ den Werth des Integrals
\[\int e^{-x}x^z dx\]
\noindent von $x=0$ bis $x=\infty$ bezeichnet, welcher, falls $z$ eine ganze Zahl ist, durch das Product $1{.}\,2{.}\,3{.}\,4\,{.}\,{.}\,{.}\,z$ ausgedrückt wird.
\begin{equation}
\begin{aligned}
&F(\alpha,\beta,\gamma,1)=\frac{i^{\gamma-1}F(-i,\gamma-1,\gamma,1)}{i^{\gamma-1-\alpha}F(-i,\gamma-1-\alpha,\gamma-\alpha,1)}\mathpunct{\raisebox{.35ex}{.}}\\[.7em]
&\qquad\mathpunct{\raisebox{.35ex}{.}}\frac{i^{\gamma-1-\alpha-\beta}F(-i,\gamma-1-\alpha-\beta,\gamma-\alpha-\beta,1)}{i^{\gamma-\beta-1}F(-i,\gamma-1-\beta,\gamma-\beta,1)}F(\alpha,\beta,i+\gamma,1),
\end{aligned}\tag*{IV.}
\end{equation}
\begin{equation}
\begin{aligned}
F\left(\alpha,\beta,\tfrac12,x\right)&=\mathfrak{A}F\left(2\alpha,2\beta,\alpha+\beta+\tfrac12,\frac{1+\surd x}{2}\right)\\[.7em]
&\qquad+\mathfrak{A}F\left(2\alpha,2\beta,\alpha+\beta+\tfrac12,\frac{1-\surd x}{2}\right),
\end{aligned}\tag*{V.}
\end{equation}
\noindent wo
\[\mathfrak{A}=\frac{\Pi(-\tfrac12){.}\,\Pi(-\tfrac12-\alpha+\beta){.}\,\Pi(-\tfrac12+\alpha-\beta)}{\Pi(-\tfrac12-\alpha){.}\,\Pi(-\tfrac12-\beta){.}\,\Pi(-\tfrac12-\alpha-\beta)};\]
\begin{equation}
\begin{aligned}
F\left(\alpha,\beta,\tfrac32,x\right)&=\frac{\mathfrak{B}}{\surd x}F\left(2\alpha-1,2\beta-1,\alpha+\beta-\tfrac12,\frac{1+\surd x}{2}\right)\\[.7em]
&\qquad-\frac{\mathfrak{B}}{\surd x}F\left(2\alpha-1,2\beta-1,\alpha+\beta-\tfrac12,\frac{1-\surd x}{2}\right),
\end{aligned}\tag*{VI.}
\end{equation}
\clearpage\SourceLeaf{359}\ArithmeticHead{352}{L}{ANALYSIS.\quad NACHLASS.}
\noindent wo
\[\mathfrak{B}=\frac{2\alpha+2\beta-1}{(2\alpha-1)(2\beta-1)}\mathfrak{A}.\]
\clearpage
\endgroup
